% GENERATED FILE. Edit canonical lessons, then run npm run generate:books.
% canonical-source-hash: fnv1a64:edbe1863e8256423
% canonical-lessons: FA-C359-zafaran, FA-C359-ezafe-kardan, FA-C359-saf-kardan, FA-C359-bereshte, FA-C359-tam, FA-R359-review-356-359, FA-R359-review-final-359

\chapter{Finishing the dish}
\label{ch:fa-finishing-the-dish}
\hlchaptermodality{359}

\begin{chapteropening}
\textbf{By the end of this chapter:} \emph{I can talk about saffron, adding something, straining, a crisp golden finish and the taste of a dish.}

Complete the last lesson of chapter 359: I can talk about saffron, adding something, straining, a crisp golden finish and the taste of a dish.
\end{chapteropening}

\section[\texorpdfstring{\fa{زعفران} (za'farân)}{زعفران (za'farân)}]{\fa{زعفران} (za'farân) — saffron}

\label{lesson:FA-C359-zafaran}



\begin{quote}
Before the new one: say the Persian for to be well cooked, then the Persian for a slotted spoon.
\end{quote}

\subsection*{You'll want to know: \fa{زعفران}}
\textbf{\fa{زعفران}} (\emph{za'farân}) — \textquotedblleft{}saffron\textquotedblright{}.

\begin{quote}
\fa{زعفران} \fa{به} \fa{برنج} \fa{رنگ} \fa{زرد} \fa{میدهد}.
\end{quote}

\begin{quote}
\emph{Za'farân be berenj rang-e zard midahad.}
\end{quote}

\begin{quote}
Saffron gives rice a yellow colour.
\end{quote}

It is brewed first in a little hot water, then poured over the rice.

\begin{cousinweb}
\emph{za'farân} is borrowed from Arabic \emph{za'farān}, \textquotedblleft{}saffron\textquotedblright{} (EtymDB), and EtymDB links the Arabic word in turn with a Persian \emph{zarparān}, \textquotedblleft{}with golden leaves\textquotedblright{}.
\end{cousinweb}

\begin{grammarlens}[title={What you've built}]
One more word for telling, explaining and describing at B1.
\end{grammarlens}

\subsection*{Your turn}
\begin{itemize}
\raggedright
  \item \emph{Say it:} \emph{za'farân}
  \item \emph{Say it:} \emph{za'farân}, once more
  \item \emph{Say it:} a sentence of your own with \emph{za'farân}
  \item \emph{Recall:} say the Persian for to be well cooked, then the Persian for a slotted spoon, then say \emph{za'farân} again
\end{itemize}

\begin{tcolorbox}[breakable,colback=teal!4,colframe=teal!35!black,title={Before you move on}]
What does \fa{زعفران} mean? (\textquotedblleft{}Saffron\textquotedblright{}.) Say it once more.
\end{tcolorbox}
\section[\texorpdfstring{\fa{اضافه} \fa{کردن} (ezâfe kardan)}{اضافه کردن (ezâfe kardan)}]{\fa{اضافه} \fa{کردن} (ezâfe kardan) — to add}

\label{lesson:FA-C359-ezafe-kardan}



\begin{quote}
Before the new one: say the Persian for a slotted spoon, then the Persian for saffron.
\end{quote}

\subsection*{You'll want to know: \fa{اضافه} \fa{کردن}}
\textbf{\fa{اضافه} \fa{کردن}} (\emph{ezâfe kardan}) — \textquotedblleft{}to add\textquotedblright{}.

\begin{quote}
\fa{در} \fa{آخر} \fa{کمی} \fa{نمک} \fa{اضافه} \fa{کرد}.
\end{quote}

\begin{quote}
\emph{Dar âkhar kami namak ezâfe kard.}
\end{quote}

\begin{quote}
At the end she added a little salt.
\end{quote}

With \emph{be}: \emph{be khoresh ezâfe kard}. \emph{Ezâfe} is also the name of the Persian linking \emph{-e}.

\begin{cousinweb}
\emph{ezâfe} is borrowed from Arabic \emph{iḍāfa}, \textquotedblleft{}an adding, an adjunct\textquotedblright{} (EtymDB). \emph{kardan} \textquotedblleft{}to do\textquotedblright{} is inherited Persian.
\end{cousinweb}

\begin{grammarlens}[title={What you've built}]
One more word for telling, explaining and describing at B1.
\end{grammarlens}

\subsection*{Your turn}
\begin{itemize}
\raggedright
  \item \emph{Say it:} \emph{ezâfe kardan}
  \item \emph{Say it:} \emph{ezâfe kardan}, once more
  \item \emph{Say it:} a sentence of your own with \emph{ezâfe kardan}
  \item \emph{Recall:} say the Persian for a slotted spoon, then the Persian for saffron, then say \emph{ezâfe kardan} again
\end{itemize}

\begin{tcolorbox}[breakable,colback=teal!4,colframe=teal!35!black,title={Before you move on}]
What does \fa{اضافه} \fa{کردن} mean? (\textquotedblleft{}To add\textquotedblright{}.) Say it once more.
\end{tcolorbox}
\section[\texorpdfstring{\fa{صاف} \fa{کردن} (sâf kardan)}{صاف کردن (sâf kardan)}]{\fa{صاف} \fa{کردن} (sâf kardan) — to strain, to filter; to smooth}

\label{lesson:FA-C359-saf-kardan}



\begin{quote}
Before the new one: say the Persian for saffron, then the Persian for to add.
\end{quote}

\subsection*{You'll want to know: \fa{صاف} \fa{کردن}}
\textbf{\fa{صاف} \fa{کردن}} (\emph{sâf kardan}) — \textquotedblleft{}to strain, to filter; to smooth\textquotedblright{}.

\begin{quote}
\fa{چای} \fa{را} \fa{صاف} \fa{کرد} \fa{و} \fa{در} \fa{لیوان} \fa{ریخت}.
\end{quote}

\begin{quote}
\emph{Chây râ sâf kard va dar livân rikht.}
\end{quote}

\begin{quote}
She strained the tea and poured it into a glass.
\end{quote}

Also \textquotedblleft{}to smooth, to straighten\textquotedblright{}. Present \emph{sâf mikonad}.

\begin{cousinweb}
\emph{sâf} \textquotedblleft{}clear, smooth\textquotedblright{} is borrowed from Arabic \emph{ṣāfī}, \textquotedblleft{}pure\textquotedblright{} (EtymDB). \emph{kardan} \textquotedblleft{}to do\textquotedblright{} is inherited Persian: to make clear.
\end{cousinweb}

\begin{grammarlens}[title={What you've built}]
One more word for telling, explaining and describing at B1.
\end{grammarlens}

\subsection*{Your turn}
\begin{itemize}
\raggedright
  \item \emph{Say it:} \emph{sâf kardan}
  \item \emph{Say it:} \emph{sâf kardan}, once more
  \item \emph{Say it:} a sentence of your own with \emph{sâf kardan}
  \item \emph{Recall:} say the Persian for saffron, then the Persian for to add, then say \emph{sâf kardan} again
\end{itemize}

\begin{tcolorbox}[breakable,colback=teal!4,colframe=teal!35!black,title={Before you move on}]
What does \fa{صاف} \fa{کردن} mean? (\textquotedblleft{}To strain, to filter; to smooth\textquotedblright{}.) Say it once more.
\end{tcolorbox}
\section[\texorpdfstring{\fa{برشته} (bereshte)}{برشته (bereshte)}]{\fa{برشته} (bereshte) — toasted, crisp and golden}

\label{lesson:FA-C359-bereshte}



\begin{quote}
Before the new one: say the Persian for to add, then the Persian for to strain.
\end{quote}

\subsection*{You'll want to know: \fa{برشته}}
\textbf{\fa{برشته}} (\emph{bereshte}) — \textquotedblleft{}toasted, crisp and golden\textquotedblright{}.

\begin{quote}
\fa{نان} \fa{برشته} \fa{بوی} \fa{خوبی} \fa{داشت}.
\end{quote}

\begin{quote}
\emph{Nân-e bereshte bu-ye khubi dâsht.}
\end{quote}

\begin{quote}
The toasted bread smelled good.
\end{quote}

From \emph{bereshtan} \textquotedblleft{}to roast\textquotedblright{}. \emph{Qahve-ye bereshte} \textquotedblleft{}roasted coffee\textquotedblright{}.

\begin{cousinweb}
\emph{bereshte} is the participle of the native verb \emph{bereshtan} \textquotedblleft{}to roast, to parch\textquotedblright{}; EtymDB lists no entry.
\end{cousinweb}

\begin{grammarlens}[title={What you've built}]
One more word for telling, explaining and describing at B1.
\end{grammarlens}

\subsection*{Your turn}
\begin{itemize}
\raggedright
  \item \emph{Say it:} \emph{bereshte}
  \item \emph{Say it:} \emph{bereshte}, once more
  \item \emph{Say it:} a sentence of your own with \emph{bereshte}
  \item \emph{Recall:} say the Persian for to add, then the Persian for to strain, then say \emph{bereshte} again
\end{itemize}

\begin{tcolorbox}[breakable,colback=teal!4,colframe=teal!35!black,title={Before you move on}]
What does \fa{برشته} mean? (\textquotedblleft{}Toasted, crisp and golden\textquotedblright{}.) Say it once more.
\end{tcolorbox}
\section[\texorpdfstring{\fa{طعم} (ta'm)}{طعم (ta'm)}]{\fa{طعم} (ta'm) — a taste, a flavour}

\label{lesson:FA-C359-tam}



\begin{quote}
Before the new one: say the Persian for to strain, then the Persian for toasted.
\end{quote}

\subsection*{You'll want to know: \fa{طعم}}
\textbf{\fa{طعم}} (\emph{ta'm}) — \textquotedblleft{}a taste, a flavour\textquotedblright{}.

\begin{quote}
\fa{طعم} \fa{این} \fa{خورش} \fa{خیلی} \fa{خوب} \fa{است}.
\end{quote}

\begin{quote}
\emph{Ta'm-e in khoresh kheyli khub ast.}
\end{quote}

\begin{quote}
This stew tastes very good.
\end{quote}

\emph{Khoshta'm} \textquotedblleft{}tasty\textquotedblright{}. Close to \emph{maze}.

\begin{cousinweb}
\emph{ta'm} is borrowed from Arabic \emph{ṭa'm}, \textquotedblleft{}taste\textquotedblright{} (EtymDB).
\end{cousinweb}

\begin{grammarlens}[title={What you've built}]
One more word for telling, explaining and describing at B1.
\end{grammarlens}

\subsection*{Your turn}
\begin{itemize}
\raggedright
  \item \emph{Say it:} \emph{ta'm}
  \item \emph{Say it:} \emph{ta'm}, once more
  \item \emph{Say it:} a sentence of your own with \emph{ta'm}
  \item \emph{Recall:} say the Persian for to strain, then the Persian for toasted, then say \emph{ta'm} again
\end{itemize}

\begin{tcolorbox}[breakable,colback=teal!4,colframe=teal!35!black,title={Before you move on}]
What does \fa{طعم} mean? (\textquotedblleft{}A taste, a flavour\textquotedblright{}.) Say it once more.
\end{tcolorbox}
\section[(dialogue)]{Step by step, once more}

\label{lesson:FA-R359-review-356-359}



\begin{quote}
Say the Persian for to chop finely and to mince.
\end{quote}

\subsection*{Your turn}
Say these aloud:

\begin{itemize}
\raggedright
  \item to chop finely, to mince, to mash, to sauté, fried onion
  \item dough, to knead, to rise, yeast, a clay oven
  \item a cooking pot, to come to the boil, to bubble, to be well cooked, a slotted spoon
  \item saffron, to add, to strain, toasted, a taste
\end{itemize}

\begin{tcolorbox}[breakable,colback=teal!4,colframe=teal!35!black,title={Before you move on}]
To chop finely? (\textbf{\fa{ریز} \fa{کردن}}.) To mince? (\textbf{\fa{چرخ} \fa{کردن}}.) To mash? (\textbf{\fa{له} \fa{کردن}}.) To sauté? (\textbf{\fa{تفت} \fa{دادن}}.) Fried onion? (\textbf{\fa{پیاز} \fa{داغ}}.) Dough? (\textbf{\fa{خمیر}}.) To knead? (\textbf{\fa{ورز} \fa{دادن}}.) To rise? (\textbf{\fa{ور} \fa{آمدن}}.) Yeast? (\textbf{\fa{خمیرمایه}}.) A clay oven? (\textbf{\fa{تنور}}.) A cooking pot? (\textbf{\fa{قابلمه}}.) To come to the boil? (\textbf{\fa{جوش} \fa{آمدن}}.) To bubble? (\textbf{\fa{قل} \fa{زدن}}.) To be well cooked? (\textbf{\fa{جا} \fa{افتادن}}.) A slotted spoon? (\textbf{\fa{کفگیر}}.) Saffron? (\textbf{\fa{زعفران}}.) To add? (\textbf{\fa{اضافه} \fa{کردن}}.) To strain? (\textbf{\fa{صاف} \fa{کردن}}.) Toasted? (\textbf{\fa{برشته}}.) A taste? (\textbf{\fa{طعم}}.)
\end{tcolorbox}
\section[(dialogue)]{Finishing the dish — once more}

\label{lesson:FA-R359-review-final-359}



\begin{quote}
Say the Persian for saffron and to add.
\end{quote}

\subsection*{Your turn}
\begin{itemize}
\raggedright
  \item \emph{Say it:} saffron, to add, to strain, toasted, a taste
\end{itemize}

\begin{tcolorbox}[breakable,colback=teal!4,colframe=teal!35!black,title={Before you move on}]
Saffron? (\textbf{\fa{زعفران}}.) To add? (\textbf{\fa{اضافه} \fa{کردن}}.) To strain? (\textbf{\fa{صاف} \fa{کردن}}.) Toasted? (\textbf{\fa{برشته}}.) A taste? (\textbf{\fa{طعم}}.)
\end{tcolorbox}
